\documentclass[11pt]{article}
\usepackage[T1]{fontenc}
\usepackage{lmodern}
\usepackage[margin=1in]{geometry}
\usepackage{amsmath,amssymb,amsthm,booktabs}
\usepackage[colorlinks=true,urlcolor=blue,linkcolor=blue,citecolor=blue]{hyperref}
\usepackage{microtype}
\hypersetup{pdftitle={A Mellin-averaged determinant bound for the tenth discriminant},
 pdfauthor={Research note prepared with Codex},
 pdfsubject={Version 0.5.0: q < 2000}}
\newtheorem{theorem}{Theorem}
\newtheorem{lemma}[theorem]{Lemma}
\newtheorem{proposition}[theorem]{Proposition}
\newcommand{\Z}{\mathbb Z}
\newcommand{\Q}{\mathbb Q}
\newcommand{\C}{\mathbb C}
\newcommand{\OK}{\mathcal O_K}
\newcommand{\Tr}{\operatorname{Tr}}
\newcommand{\Res}{\operatorname{Res}}
\newcommand{\ord}{\operatorname{ord}}
\title{A Mellin-averaged determinant bound for the tenth discriminant\\
\large Version 5: \(q<2\,000\)}
\author{Research note prepared with Codex}
\date{Version 0.5.0 --- 20 October 2026}
\begin{document}
\maketitle

\begin{abstract}
We improve the analytic cutoff for an imaginary quadratic field of class
number one from \(q<4000\) to \(q<2000\). The new step averages a derivative
determinant inequality over its size parameter, producing a lower bound
for a convergent weighted sum of non-inert primes. The Euler product of
primitive quadratic norms gives a competing upper bound that accounts
for composite norms. An elementary summation formula controls the
quadratic norm series and its derivative. Exact rational intervals give
a contradiction gap greater than \(0.017\) throughout
\(2000\le q\le4000\). The independently verified v4 paper is cited for
\(q\ge4000\); its infinite-range proof is not repeated. The argument for
the new interval, including the determinant and resultant lemmas, is
given in full. A separate finite computation leaves the usual nine
discriminants.
\end{abstract}

\paragraph{Status and attribution.}
This is an AI-assisted v5 research draft. Codex developed the new
averaging argument, estimates, exposition, and verification programs.
The v5 argument has passed a manual mathematical audit with exact
computational checks. This revision applies its wording correction in
Section 6.2; the bound and certificate are unchanged. This is not formal
verification. Green's note \cite{green} supplies the
derivative determinant, Frobenius row operation, and resultant method.
The sharpened rank criterion and lattice construction are retained from
v4 \cite{v4}; the averaging and primitive-norm comparison are the new
steps here.

\section{Result, dependencies, and proof strategy}

Let \(D<0\) be a fundamental discriminant, set \(q=|D|\), and put
\(K=\Q(\sqrt D)\). Write \(\OK\) for the ring of integers and \(h(K)\)
for its ideal class number. All logarithms are natural.

\begin{theorem}\label{main}
If \(h(K)=1\), then \(q<2000\). The finite verification supplied with
this paper consequently gives
\[
 h(K)=1\quad\Longleftrightarrow\quad
 q\in\{3,4,7,8,11,19,43,67,163\}.
\]
\end{theorem}

The only earlier-version input is the already established implication
\[
 h(K)=1\ \Longrightarrow\ q<4000,
\]
from the fixed v4 release \cite{v4}. That infinite-range argument is
referenced rather than reproduced. Everything needed for
the new exclusion of \(2000\le q\le4000\) is set out below, relative
to the stated classical inputs: quadratic ideal theory,
Rosser--Schoenfeld's ordinary-prime estimates, elementary
Euler--Maclaurin summation and gamma identities. Reduction of binary
quadratic forms is used only for the finite verification.

The proof of the new interval has three parts. First, a positive integer
Gram determinant bounds the weighted inert-prime sum for every
sufficiently large determinant parameter. Averaging this family of
inequalities gives a lower bound for a convergent sum over the other
primes. Second, unique factorization of ideals expresses the same sum
through the Euler product of primitive norms. Third, an explicit
summation estimate for the norm series gives an incompatible upper
bound. Monotonicity reduces the entire real interval
\([2000,4000]\) to one endpoint; no discriminants in that interval
are individually enumerated to establish the analytic improvement.

We use no Heegner method or computed zero-free region for quadratic
Dirichlet \(L\)-functions. The ordinary Riemann zeta function below is
evaluated only at real arguments greater than one.

\section{Elementary consequences of class number one}\label{elementary}

For the new analytic argument assume
\[
 h(K)=1,\qquad 2000\le q\le4000,\qquad u=\sqrt q.
\]
A split or ramified prime \(p\) has an integral ideal of norm \(p\).
Since the ideal is principal, its generator has absolute norm \(p\).
In a standard integral basis, a nonrational quadratic integer has
imaginary part of absolute value at least \(u/2\), and hence norm at
least \(q/4\). A rational integer has square norm. Therefore
\begin{equation}\label{small-inert}
 p<q/4\quad\Longrightarrow\quad p\text{ is inert in }K.
\end{equation}
This argument is valid for any class-number-one field, independently of
the interval restriction.

In particular, 2 is inert, so \(D\equiv5\pmod8\), or \(q\equiv3\pmod8\).
Thus \(q\) is squarefree and
\[
 \omega=\frac{1+i\sqrt q}{2},\qquad
 \OK=\Z+\Z\omega,\qquad
 Q(a,b)=|a+b\omega|^2=(a+b/2)^2+qb^2/4.
\]
Moreover, \(q\) is prime: otherwise its smallest prime divisor is at
most \(\sqrt q<q/4\), but is ramified, contradicting
\eqref{small-inert}. The only ramified prime is consequently \(q\).
The only units are \(\pm1\), since \(q>4\).

\section{Explicit estimates for ordinary primes}\label{primes}

Put \(A_{\rm pr}(x)=\sum_{p\le x}\log p/p\) and
\(\vartheta(x)=\sum_{p\le x}\log p\). Theorems 6 and 9 of
Rosser--Schoenfeld \cite{rs} give
\[
 |A_{\rm pr}(x)-\log x-c_0|<\frac1{2\log x}\quad(x\ge319),
 \qquad \vartheta(x)<cx,\quad c=1.01624,\quad x>0.
\]
The finite sum at 319 gives
\(\log319-A_{\rm pr}(319)+1/(2\log319)<1.334\), so
\begin{equation}\label{prime-total}
 A_{\rm pr}(x)>\log x-1.334-\frac1{2\log x}\quad(x\ge319).
\end{equation}
Partial summation also gives \(A_{\rm pr}(x)<2\log x\) for \(x\ge2\).
The stated inequalities for the finite sums are included in the
certificate.

We use the fixed prime cutoff \(H=110\), and write
\[
 a_H=\sum_{p\le H}\frac{\log p}{p},\quad
 \theta_H=\sum_{p\le H}\log p,\quad
 W=\sum_{p\le H}\left(\frac1{110}-\frac{p}{3\cdot110^2}\right)\log p.
\]
Exact logarithm intervals give
\begin{equation}\label{prime-constants}
 a_H<3.546874,\qquad \theta_H<103,\qquad W>0.759467.
\end{equation}
For the prime-square tail, partial summation at \(N=10000\) gives
\[
 \sum_{p>N}\frac{\log p}{p^2}
 <\frac{2c}{N}-\frac{\vartheta(N)}{N^2}.
\]
Adding the explicit primes between \(H\) and \(N\), the certificate
therefore proves
\begin{equation}\label{prime-tail}
 \sum_{p>110}\frac{\log p}{p^2}<b:=0.00820476.
\end{equation}
The inertness of every prime at most 110 follows from
\eqref{small-inert}.

\section{A sharpened derivative interpolation lemma}\label{rank}

This section restates the v4 refinement of Green's resultant proof \cite{green,v4}. We retain the
actual constant in the bound for the number of points on a curve, and
retain the loss of \(w\)-degree when removing factors depending only on
\(w\). These changes give the following useful criterion.

\begin{lemma}\label{rank-lemma}
Let \(P_i=(z_i,w_i)\in\C^2\), \(1\le i\le n\), with all \(w_i\)
distinct. Suppose every nonzero polynomial \(F\in\C[z,w]\) satisfies
\[
 \#\{i:F(P_i)=0\}\le A\deg F.
\]
Let \(R,S,J\) be positive integers with
\begin{equation}\label{rank-condition}
 J\le2R+A,\qquad nJ>2RS+A(R+S).
\end{equation}
Then the \(RS\) vectors
\[
 v_{r,s}=\bigl(\partial_z^j(z^r w^s)(P_i)\bigr)_{
                 1\le i\le n,\ 0\le j<J},
 \qquad0\le r<R,\quad0\le s<S,
\]
are linearly independent over \(\C\).
\end{lemma}

\begin{proof}
A dependence would give a nonzero polynomial of degrees less than
\((R,S)\) whose first \(J\) derivatives in \(z\) vanish at every point.
Factor it over \(\C\) as
\[
 F(z,w)=F_0(w)\prod_{\ell=1}^k F_\ell(z,w),
\]
where each \(F_\ell\) is irreducible and has positive \(z\)-degree
\(d_\ell\); factors may be repeated. Write \(d'_\ell=\deg_w F_\ell\)
and \(d'_0=\deg F_0\). We can use the weak bounds
\(\sum d_\ell\le R\) and \(\sum d'_\ell\le S-d'_0\).

For an irreducible factor of bidegree \((d,d')\) with \(d>0\), the
polynomial \(T(w)=\Res_z(F_\ell,\partial_zF_\ell)\) is nonzero and
has degree at most \(2dd'\). No specialization
\(F_\ell(z,w_i)\) is identically zero: otherwise \(w-w_i\) would
divide the irreducible polynomial. At a root of multiplicity \(m\),
\begin{equation}\label{resultant-order}
 m\le1+\ord_{w=w_i}T(w).
\end{equation}
Here is a direct justification. The Sylvester matrix has size \(2d-1\).
After specialization at \(w_i\), both polynomials are divisible by
\((z-z_i)^{m-1}\), so the matrix has rank deficiency at least \(m-1\).
A constant invertible row transformation can make that many rows zero
at \(w_i\). Each such row is then divisible by \(w-w_i\), so the
determinant vanishes to order at least \(m-1\). Its entries have
\(w\)-degree at most \(d'\), which also proves the degree bound.
Nonvanishing follows from irreducibility and characteristic zero.

The number of points where this factor vanishes is at most
\(A(d+d')\). Summing \eqref{resultant-order} over these points, using
distinct \(w_i\), bounds the sum of its specialized multiplicities by
\(A(d+d')+2dd'\).

There are at most \(d'_0\) indices where \(F_0(w_i)=0\). At every
other index the specialization of \(F\) has multiplicity at least \(J\).
Thus
\begin{align*}
 J(n-d'_0)
 &\le A(R+S-d'_0)+2R(S-d'_0),\\
 nJ
 &\le A(R+S)+2RS+(J-A-2R)d'_0
 \le A(R+S)+2RS.
\end{align*}
The last inequality uses \(J\le2R+A\), and contradicts
\eqref{rank-condition}.
\end{proof}

\section{The integral determinant and its prime divisibility}\label{det}

Take distinct points \(z_i\in\OK\) and \(w_i=\overline{z_i}\) satisfying
Lemma~\ref{rank-lemma}. The derivative vectors lie in \(\OK^{nJ}\).
Form the real Gram matrix
\[
 G_{(r,s),(r',s')}=
 \Tr_{K/\Q}\!\left(\sum_{i,j}
 (v_{r,s})_{i,j}\,\overline{(v_{r',s'})_{i,j}}\right).
\]
Every entry is an integer. The trace pairing on this imaginary
quadratic field is the positive real inner product
\(2\operatorname{Re}(\alpha\overline\beta)\). Complex linear
independence of the vectors implies their real linear independence.
Consequently
\begin{equation}\label{positive-det}
 \Delta:=\det G\in\Z_{>0}.
\end{equation}

\subsection{The inert-prime bound, including the prime 2}

If \(p\) is inert, \(\OK/p\OK\) is the field of \(p^2\) elements.
Its Frobenius is the nontrivial quadratic automorphism, so
\(z_i^p\equiv\overline{z_i}=w_i\pmod{p\OK}\).
This applies to \(p=2\) as well as odd inert primes because we work
in the full ring of integers. Put \(\psi=z^p-w\).
Its value and all its positive \(z\)-derivatives at these points are
in \(p\OK\): every nonzero derivative coefficient contains \(p\).

For each pair \((r,s)\), let
\(h_{r,s}=\min(\lfloor r/p\rfloor,S-1-s)\).
Replace its polynomial by
\(z^{r-ph_{r,s}}\psi^{h_{r,s}}w^s\).
The binomial expansion has leading term \(z^r w^s\) and other terms
\(z^{r-pi}w^{s+i}\), all with smaller \(r\) and still within the
rectangle of indices. The induced row transformation is an integer
triangular matrix with determinant one. Every derivative of the
replacement polynomial is in \(p^{h_{r,s}}\OK\), by the product rule.
Applying the same transformation to the Gram matrix and factoring
from rows and columns proves
\begin{equation}\label{valuation}
 v_p(\Delta)\ge
 2\sum_{0\le r<R,\ 0\le s<S}\min(\lfloor r/p\rfloor,S-1-s).
\end{equation}
No division by derivative factorials is made in this argument.

\subsection{Explicit lower bounds for the valuation sum}

Let \(e(R,S,p)\) denote the right side of \eqref{valuation}, divided
by \(R^2S\). Counting the pairs for which the minimum is at least
\(h\) gives the exact identity
\begin{equation}\label{layers}
 e(R,S,p)=\frac2{R^2S}
 \sum_{h=1}^{S-1}(S-h)\max(R-ph,0).
\end{equation}
When \(p(S-1)\le R-1\), summing this quadratic expression gives
\begin{equation}\label{small-e}
 e(R,S,p)=\frac{S-1}{R}-\frac{p(S^2-1)}{3R^2}.
\end{equation}
For a bound valid for all \(p\), use
\(\sum_{s=0}^{S-1}\min(k,S-1-s)\ge Sk-k(k+1)/2\).
Writing \(R=ap+b\), \(0\le b<p\), shows
\[
 \sum_{r=0}^{R-1}\lfloor r/p\rfloor
 =\frac{R(R-p)}{2p}+\frac{b(p-b)}{2p}
 \ge\frac{R(R-p)}{2p}.
\]
Also \(\lfloor r/p\rfloor(\lfloor r/p\rfloor+1)
\le r^2/p^2+r/p\). Substitution gives
\begin{equation}\label{tail-e}
 e(R,S,p)\ge
 \left(1-\frac{R-1}{2RS}\right)\frac1p
 -\frac1R-\frac{(R-1)(2R-1)}{6RS\,p^2}.
\end{equation}

\section{A determinant inequality at every large scale}\label{scales}

Set
\[
 t=\frac1{110},\quad j_0=\frac1{160},\quad
 C_0=\frac{16}{5},\quad \eta=\frac{2001}{2000},\quad
 M=50\,000\,000.
\]
For each integer \(R\ge M\), choose
\[
 S=\lfloor tR\rfloor,\qquad J=\lfloor j_0R\rfloor,\qquad
 N_x=\left\lceil\sqrt{C_0Ru/2}\right\rceil,\quad
 N_y=\left\lceil\sqrt{2C_0R/u}\right\rceil.
\]
Take \(N_y\) consecutive integers \(b\), centered so that
\(|b|\le N_y/2\). For each \(b\), take \(N_x\) consecutive integers
\(a\) with \(|a+b/2|\le N_x/2\), and use the points
\((z,w)=(a+b\omega,a+b\overline\omega)\).
There are \(n=N_xN_y\) points and their \(w\)-coordinates are distinct.

Both unrounded side lengths exceed 2000, since \(44<u<64\).
Consequently
\[
 C_0R\le n<\eta^2C_0R<4R,\qquad
 |z|=|w|\le T:=\eta\sqrt{C_0Ru/4},\qquad T/R<0.00102.
\]
Here \(T\ge1\). The zero-count constant in
Lemma~\ref{rank-lemma} can be taken to be \(A=N_x\).
Indeed, after the invertible linear change from \((z,w)\) to \((a,b)\),
a polynomial of total degree \(d\) that vanishes identically on \(k\)
horizontal lines has at most
\[
 kN_x+(N_y-k)(d-k)\le N_xd
\]
zeros on the grid. This follows by factoring out those \(k\) line
factors; \(N_x\ge N_y\) and \(k\le d\).

Uniformly in the stated ranges, \(A/R<0.00144\), and
\[
 C_0j_0-2t>
 \frac{C_0}{M}+0.00144(1+t).
\]
It follows that \(nJ>2RS+A(R+S)\).
The other rank condition \(J\le2R+A\) is immediate.
Thus the positive integer determinant of Section~\ref{det} exists
for every integer \(R\ge M\).

\subsection{Averaging the derivative size over the rows}

There is a small improvement to the usual maximum-row estimate.
For a monomial \(z^rw^s\), its derivative entries satisfy
\[
 |\partial_z^k(z^rw^s)|\le
 T^{r+s}\max(1,r/T)^{J-1}\qquad(0\le k<J).
\]
The derivative is zero if \(k>r\), which is consistent with this bound.
Since the function \(\log_+(x/T)=\max(0,\log(x/T))\) is increasing,
\[
 \sum_{r=0}^{R-1}\log_+(r/T)
 \le\int_0^R\log_+(x/T)\,dx=R\log(R/T)-R+T.
\]
Hadamard's inequality now gives
\begin{equation}\label{mean-arch}
 \frac{\log\Delta}{R^2S}
 \le \frac{\log(2nJ)}R+
 \left(1+\frac SR\right)\log T+
 \frac{2J}R\bigl(\log(R/T)-1+T/R\bigr).
\end{equation}
The last bracket is nonnegative, so \(S/R\le t\) and \(J/R\le j_0\)
can be substituted on the right. Put
\[
 \beta=1+t-2j_0=\frac{877}{880},\qquad
 \alpha=1-\frac{1+t}{2}-j_0=\frac{861}{1760}.
\]
The prefactor \(\log(2nJ)/R\) is less than \(10^{-6}\), and
\(2j_0T/R+10^{-6}<0.00002\). Therefore
\begin{equation}\label{new-arch}
 \frac{\log\Delta}{R^2S}
 <(1-\alpha)\log R+\frac{\beta}{4}\log q+
 \frac{\beta}{2}\log(4/5)+\beta\log\eta-2j_0+0.00002.
\end{equation}

\subsection{The prime divisibility comparison}

Write
\[
 I_R=\sum_{\substack{110<p\le R\\p\ {\rm inert}}}\frac{\log p}{p},
 \qquad
 S_-(R)=\sum_{\substack{p\le R\\p\ {\rm inert}}}\frac{\log p}{p}.
\]
The small-prime valuation formula \eqref{small-e} gives a total
contribution at least \(W-2\theta_H/R\); its hypothesis holds since
the largest prime at most 110 is 109.
For the other inert primes, \eqref{tail-e} gives
\[
 \frac{\log\Delta}{R^2S}>
 I_R+W-c-\frac{110.001}{3}b-0.0001.
\]
In detail, \(R/S\le110.001\), the small-prime floor loss is less than
\(0.000005\), and
\[
 \frac{I_R}{2S}<\frac{\log R}{S}
 \le110.001\frac{\log R}{R}
 \le110.001\frac{\log M}{M}<0.00004.
\]
Here \((\log R)/R\) decreases for \(R\ge M>e\).
The final inequality is the one checked by the certificate; no
monotonicity of \((\log R)/S\) is needed.
The loss from \(\vartheta(R)/R\) is at most \(c\).
These estimates explain the stated \(0.0001\) allowance.

Combining this with \eqref{new-arch}, and adding \(a_H\), yields
\begin{equation}\label{inert-all-scales}
 S_-(x)<(1-\alpha)\log x+\frac{\beta}{4}\log q+C_*,
 \qquad C_*=3.982,\qquad x\ge M.
\end{equation}
The certificate checks
\[
 a_H-W+c+\frac{110.001}{3}b+
 \frac{\beta}{2}\log(4/5)+\beta\log\eta-2j_0+0.00012<3.982.
\]
Initially the argument proves \eqref{inert-all-scales} for integer \(x\).
For real \(x\ge M\), apply it to \(R=\lfloor x\rfloor\): the prime
sum is unchanged and \(\log R\le\log x\).

\section{Mellin averaging gives a lower bound}\label{mellin}

Fix
\[
 \delta=\frac1{70},\qquad \sigma=1+\delta=\frac{71}{70},
 \qquad m=\log M.
\]
Let \(S_{\ne-}(x)\) be the sum of \(\log p/p\) over split and ramified
primes at most \(x\). Equations \eqref{prime-total} and
\eqref{inert-all-scales} imply
\[
 S_{\ne-}(x)>
 \alpha\log x-\frac{\beta}{4}\log q-C_*-1.334-\frac1{2m}
 \qquad(x\ge M).
\]
All terms in \(S_{\ne-}\) are nonnegative. Thus, by integrating the
step function and interchanging nonnegative sums and integrals,
\begin{align}
 D_{\ne-}(\sigma)
 &:=\sum_{p\ {\rm split\ or\ ramified}}\frac{\log p}{p^\sigma}
 =\delta\int_1^\infty S_{\ne-}(x)x^{-1-\delta}\,dx \nonumber\\
 &>P(q):=e^{-\delta m}\left[
 \alpha(m+\delta^{-1})-\frac{\beta}{4}\log q
 -3.982-1.334-\frac1{2m}\right].\label{mellin-lower}
\end{align}
This is the first new comparison: a family of determinant inequalities
produces a bound for a convergent Dirichlet sum rather than for a
single sharply truncated sum.

\section{The Euler product of primitive norms}\label{euler}

For real \(s>1\), define the nonrational norm series
\[
 A_q(s)=\sum_{b=1}^{\infty}\sum_{a\in\Z}Q(a,b)^{-s}.
\]
Because every ideal is principal and the units are \(\pm1\), the
Dedekind zeta series is
\[
 Z_K(s)=\zeta(2s)+A_q(s).
\]
The rational generators contribute \(\zeta(2s)\); each nonrational
ideal has exactly one generator with \(b>0\).

Dividing out rational integer factors gives
\begin{equation}\label{primitive-euler}
 B_q(s):=\frac{Z_K(s)}{\zeta(2s)}
 =\prod_{p\ {\rm split}}\frac{1+p^{-s}}{1-p^{-s}}\,(1+q^{-s}).
\end{equation}
One can also check this directly from ideal factorization. At a split
prime, removing a rational factor leaves a power of at most one of the
two prime ideals, giving \(1+2\sum_{k\ge1}p^{-ks}\).
At an inert prime, the sole prime ideal is generated by the rational
integer \(p\), so its primitive factor is 1. At the ramified prime \(q\),
the two possible primitive exponents are 0 and 1.
All products and differentiated series converge absolutely for \(s>1\).

Logarithmic differentiation gives
\[
 -\frac{B_q'(s)}{B_q(s)}
 =2\sum_{p\ {\rm split}}\frac{\log p}{p^s(1-p^{-2s})}
   +\frac{\log q}{q^s+1}.
\]
In particular,
\begin{equation}\label{euler-upper-start}
 D_{\ne-}(s)\le-\frac{B_q'(s)}{2B_q(s)}
                         +\frac{\log q}{q^s}.
\end{equation}
The denominator \(B_q(s)\) retains the contribution of composite
primitive norms, which was discarded when the earlier argument simply
bounded the number of prime norms by the number of primitive points.

Put \(z=\zeta(2s)\) and \(c_s=-2\zeta'(2s)/\zeta(2s)>0\).
The useful exact identity is
\begin{equation}\label{logderivative}
 -\frac{B_q'(s)}{B_q(s)}
 =\frac{-A_q'(s)-c_s A_q(s)}{z+A_q(s)}.
\end{equation}
The cancellation of the rational-generator contribution is explicit
in this formula.

\section{An explicit estimate for the norm series}\label{norm-series}

Write \(v=\sqrt q/2\), and define
\[
 C(s)=\int_{\mathbb R}(1+x^2)^{-s}\,dx
 =\frac{\sqrt\pi\,\Gamma(s-\tfrac12)}{\Gamma(s)},\qquad
 F_q(s)=C(s)v^{1-2s}\zeta(2s-1).
\]
Replacing each horizontal lattice sum by its integral gives \(F_q(s)\).
We bound the error for the series and its derivative separately; we
do not differentiate an already bounded error term.

For a twice differentiable decaying function \(f\) and a real shift
\(\tau\), Euler summation gives
\begin{equation}\label{summation-error}
 \left|\sum_{a\in\Z}f(a+\tau)-\int_{\mathbb R}f(x)\,dx\right|
 \le\frac1{12}\int_{\mathbb R}|f''(x)|\,dx.
\end{equation}
Indeed the error is, up to sign, the integral of
\(\widetilde B_2(x-\tau)f''(x)/2\), where the periodic Bernoulli
polynomial satisfies \(|\widetilde B_2|\le1/6\).
This identity follows by two integrations by parts on the unit
intervals; the boundary terms telescope. Decay justifies passage to
the whole line.

Apply this first to \(f(x)=(x^2+w^2)^{-\sigma}\) and then to
\(g(x)=\log(x^2+w^2)(x^2+w^2)^{-\sigma}\), with \(w=vb\).
Since \(\sigma<1.02\) and \(w\ge22\), elementary differentiation gives
\[
 |f''(x)|\le11(x^2+w^2)^{-\sigma-1},\qquad
 |g''(x)|\le\bigl(19+11\log(x^2+w^2)\bigr)
                         (x^2+w^2)^{-\sigma-1}.
\]
For the constants, the bounds
\(4\sigma^2+6\sigma<11\) and \(10+8\sigma<19\) suffice.
After \(x=wy\), use
\[
 \int_{\mathbb R}(1+y^2)^{-2}\,dy=\frac\pi2,\qquad
 \int_{\mathbb R}\frac{\log(1+y^2)}{(1+y^2)^2}\,dy\le\pi.
\]
The second inequality follows from \(\log t\le t\).
Summing \eqref{summation-error} over \(b\ge1\) therefore gives
\begin{align}
 |A_q(\sigma)-F_q(\sigma)|
 &\le \frac{11\pi}{24\cdot22^3}\zeta(3)
 <E_A:=0.0002,\label{norm-error}\\
 |-A_q'(\sigma)+F_q'(\sigma)|
 &\le\frac{\pi}{24\cdot22^3}
 \left((41+22\log22)\zeta(3)+22(-\zeta'(3))\right)
 <E_N:=0.002.\label{derivative-error}
\end{align}
Here \(v^{-3}(41+22\log v)\) decreases for \(v\ge22\).
Absolute convergence justifies all differentiations and summations.

From now on all functions of \(s\) are evaluated at \(s=\sigma\).
Let
\[
 \ell=\log(q/4),\quad F(q)=F_q(\sigma),\quad
 g=-\frac{C'(\sigma)}{C(\sigma)},\quad
 \kappa=g-\frac{2\zeta'(2\sigma-1)}{\zeta(2\sigma-1)}-c_\sigma .
\]
Then
\[
 -F_q'(\sigma)-c_\sigma F(q)=F(q)(\ell+\kappa).
\]
For fixed \(N\ge0,c>0,z>0\), the function
\((N-cA)/(z+A)\) decreases with \(A\). Applying this observation
to \eqref{logderivative}, with \eqref{norm-error} and
\eqref{derivative-error}, proves
\[
 -\frac{B_q'(\sigma)}{B_q(\sigma)}
 \le\frac{F(q)(\ell+\kappa)+E_N+c_\sigma E_A}
           {\zeta(2\sigma)+F(q)-E_A}.
\]
The rational certificate proves the following rounded bounds:
\[
 \zeta(2\sigma)>1.61893624,\quad
 c_\sigma<1.090928,\quad
 65<\kappa<69.106001.
\]
Thus \eqref{euler-upper-start} yields the completely specified upper
bound
\begin{equation}\label{euler-upper}
 D_{\ne-}(\sigma)\le U(q):=
 \frac{F(q)(\log(q/4)+69.106001)+0.002219}
      {2(1.61873624+F(q))}
 +\frac{\log q}{q^\sigma}.
\end{equation}

\section{One endpoint gives the new contradiction}\label{endpoint}

Throughout \(2000\le q\le4000\), the certificate proves \(3<F(q)<5\).
Set \(a=\sigma-\tfrac12>\tfrac12\),
\(A_0=1.61873624\), \(E_0=0.002219\), and \(k_0=69.106001\).
We have \(dF/d\log q=-aF\), so
\[
 \frac{d}{d\log q}\frac{F(\ell+k_0)+E_0}{A_0+F}
 =\frac{F}{A_0+F}
  \left(1-\frac{a(A_0(\ell+k_0)-E_0)}{A_0+F}\right).
\]
Use \(1.6<A_0<1.7\), \(E_0<0.003\), \(6<\ell<7\),
\(k_0>65\), and \(3<F<5\). The bracket is negative, and the
displayed derivative is at most
\[
 \frac38\left(1-\frac{(1/2)(1.6\cdot71-0.003)}{6.7}\right)<-2.
\]
The derivative of \((\log q)/q^\sigma\) with respect to \(\log q\)
is also negative. Consequently \(dU/d\log q<-1\).
In contrast, \eqref{mellin-lower} gives
\[
 \frac{dP}{d\log q}=-e^{-\delta m}\frac{\beta}{4}>-\frac14.
\]
Hence \(P(q)-U(q)\) is increasing on the whole interval.

At the single endpoint \(q=2000\), rational interval arithmetic gives
\[
 P(2000)>27.696,\qquad U(2000)<27.679.
\]
The gap is greater than \(0.017\). This contradicts
\eqref{mellin-lower} and \eqref{euler-upper} for every
\(q\in[2000,4000]\). Together with the cited v4 exclusion above
4000, it proves the analytic part of Theorem~\ref{main}.

\section{Numerical certificate and reproduction}\label{certificate}

The verification programs use only the Python standard library.
Every numerical acceptance decision uses integers or rational intervals.
The interval operations round outwards to multiples of \(10^{-45}\).
Floating-point optimization was used only to select parameters; it is
not an input to the certificate.

Logarithms are enclosed by range reduction and
\[
 \log x=2\sum_{k\ge0}\frac{t^{2k+1}}{2k+1},\qquad
 t=\frac{x-1}{x+1},\quad 1\le x\le2,
\]
using 60 terms and the positive geometric tail.
Exponentials use range reduction to \(0\le x\le1\), the Taylor
polynomial through degree 80, and remainder at most
\(x^{81}/(81!(1-x/82))\). Inversion handles negative arguments.
Machin's identity
\(\pi=16\arctan(1/5)-4\arctan(1/239)\), with 40 alternating terms
in each series, supplies the interval for \(\pi\).

For completeness, here are the special-function enclosures used in
the certificate. Write \((s)_r=s(s+1)\cdots(s+r-1)\).
Euler--Maclaurin, with \(N=64\) and \(m_0=4\), gives
\begin{align*}
 \zeta(s)={}&\sum_{n=1}^{N-1}n^{-s}+\frac{N^{1-s}}{s-1}
 +\frac12N^{-s}
 +\sum_{k=1}^{m_0}\frac{B_{2k}}{(2k)!}(s)_{2k-1}
                 N^{1-s-2k}+\mathcal R(s),\\
 |\mathcal R(s)|\le{}&
 E(s):=\frac{|B_{2m_0}|}{(2m_0)!}
       \frac{(s)_{2m_0}}{s+2m_0-1}N^{1-s-2m_0},\\
 |\mathcal R'(s)|\le{}&
 E(s)\left(\sum_{j=0}^{2m_0-1}\frac1{s+j}
             +\log N+\frac1{s+2m_0-1}\right).
\end{align*}
These estimates follow from the periodic Bernoulli remainder integral,
and from separately bounding its differentiated integrand.
The bound \(|\widetilde B_{2m_0}|\le |B_{2m_0}|\) follows from the
absolutely convergent Fourier series of that polynomial.
All values used have \(s>1\).

Finally, the beta integral and the gamma duplication formula give
\[
 C(1+\delta)=
 \pi\,2^{-2\delta}
 \frac{\Gamma(1+2\delta)}{\Gamma(1+\delta)^2}.
\]
Expanding the logarithm by Euler's gamma product yields
\begin{align*}
 \log(C(1+\delta)/\pi)
 &=-2\delta\log2+
   \sum_{n\ge2}\frac{(-1)^n(2^n-2)\zeta(n)}{n}\delta^n,\\
 -C'(1+\delta)/C(1+\delta)
 &=2\log2+
   \sum_{n\ge2}(-1)^{n+1}(2^n-2)\zeta(n)\delta^{n-1}.
\end{align*}
The linear Euler-constant terms cancel. Truncation at \(n=10\)
leaves absolute remainders at most
\[
 \frac{2(2\delta)^{11}}{1-2\delta},\qquad
 \frac{2(2\delta)^{11}}{\delta(1-2\delta)},
\]
respectively, since \(1<\zeta(n)<2\) for \(n\ge2\).
These formulas also follow from the classical expansions and
summation identities collected in \cite{dlmf}.

Run
\begin{verbatim}
python3 verify_v5_bound.py --sanity
python3 verify_finite_range.py
\end{verbatim}
The first command certifies 37 new numerical inequalities and also
runs the retained finite determinant sanity checks: 2880 valuation
cases and 60\,000 derivative row transformations. A further check
compares 10\,000 primitive-norm coefficients, enumerated directly,
with the Euler factors in five known class-number-one fields.
These finite examples do not prove the general lemmas. The second command performs
the exhaustive exclusion and reduced-form enumeration described below.

The certificate checks arithmetic and error bounds. V5 has also passed
a manual mathematical audit with exact computational checks; this revision
applies its Section 6.2 wording correction. No formal verification is claimed.

\section{The remaining finite classification}\label{finite}

For \(q>8\), the elementary norm argument in
Section 2 already forces \(q\equiv3\pmod8\) and makes every prime
\(p<q/4\) inert. Hence every remaining candidate is among the 249
integers
\[
 11\le q<2\,000,\qquad q\equiv3\pmod8.
\]
The program \texttt{verify\_finite\_range.py} checks every such integer.
For each odd prime \(p\le109\) with \(4p<q\), it rejects \(q\) if
\(-q\) is zero or a square modulo \(p\). These would respectively
make \(p\) ramified or split, contrary to \eqref{small-inert}.
Including nonfundamental integers in the scan makes the exclusion
stronger; it cannot remove a possible field incorrectly.
The surviving values are exactly
\[
 11,19,43,67,163.
\]
The negative fundamental discriminants with \(q\le8\) give
\(q=3,4,7,8\).

To check class number one for each survivor, use the correspondence
between ideal classes of discriminant \(-q\) and primitive positive
binary quadratic forms \((a,b,c)\) of discriminant
\(b^2-4ac=-q\). Each class has one reduced representative with
\[
 |b|\le a\le c,\qquad
 b\ge0\ \text{if }|b|=a\text{ or }a=c.
\]
The bound \(a\le\sqrt{q/3}\) follows from
\(q=4ac-b^2\ge3a^2\). Enumerating this finite range, retaining
\(\gcd(a,b,c)=1\), gives exactly one reduced form for each of the
nine survivors. Thus the finite program proves the remaining assertion
of Theorem~\ref{main}, conditional only on the analytic argument above
and the standard form-class correspondence.


\begin{thebibliography}{9}
\bibitem{green}
B.~Green, \emph{On Landau--Siegel zeros}, version 2, 2026.
Source snapshot identified as \texttt{2610.00311v2},
\href{https://hexagonmath.org/api/v1/works/2610.00311v2/source?path=no-siegel.tex}
{\texttt{no-siegel.tex}}.
The derivative determinant, Frobenius row operation, and resultant method
originate in this note; the needed arguments are included above.

\bibitem{rs}
J.~Barkley Rosser and L.~Schoenfeld,
\emph{Approximate formulas for some functions of prime numbers},
Illinois J. Math. \textbf{6} (1962), 64--94.
\href{https://doi.org/10.1215/ijm/1255631807}{doi:10.1215/ijm/1255631807}.
Theorems 6 and 9, printed pp.~70--71.

\bibitem{v4}
\emph{A derivative determinant bound for the tenth discriminant},
version 0.4.0, 20 October 2026.
\href{https://github.com/twinheterodox/tenth-discriminant/releases/tag/v0.4.0}
{Fixed v4 release} and
\href{https://github.com/twinheterodox/tenth-discriminant/releases/download/v0.4.0/tenth-discriminant-v0.4.0.pdf}
{direct PDF}.
The implication \(h(K)=1\Rightarrow q<4000\) is an input to this paper.
Independent mathematical verification was confirmed on 20 October 2026;
the fixed release files retain their status at publication. See the
\href{https://github.com/twinheterodox/tenth-discriminant/blob/main/REVIEW.md}
{current review record}.

\bibitem{dlmf}
NIST Digital Library of Mathematical Functions,
\href{https://dlmf.nist.gov/2.10}{Section 2.10} (Euler--Maclaurin),
\href{https://dlmf.nist.gov/5.7}{Section 5.7} (gamma expansions), and
\href{https://dlmf.nist.gov/25.2}{Section 25.2} (zeta representations).
The specific formulas and remainder bounds used in the certificate are
stated in Section~\ref{certificate}.

\bibitem{tao}
T.~Tao, \emph{An experimental AI math challenge: using tenth discriminant
bounds to explore arguments ruling out Siegel zeroes}, 20 October 2026.
\href{https://terrytao.wordpress.com/2026/10/20/an-experimental-ai-math-challenge-using-tenth-discriminant-bounds-to-explore-arguments-ruling-out-siegel-zeroes/}
{Challenge and discussion}.
\end{thebibliography}
\end{document}
